\documentclass[12pt]{article}
\usepackage{amsmath, amssymb}
\usepackage{geometry}
\usepackage{enumitem}
\geometry{margin=1in}

\begin{document}

\section*{CSE400 – Fundamentals of Probability in Computing}

\subsection*{Lecture 10: Randomized Min-Cut Algorithm}
\textbf{Instructor:} Dhaval Patel, PhD \\
\textbf{Date:} February 5, 2026

\hrule
\vspace{0.5cm}

\section*{1. Min-Cut Problem}

\subsection*{1.1 Why Use Min-Cut?}

We use the min-cut algorithm in various applications to solve problems related to \textbf{network connectivity, reliability, and optimization}.

Applications include:

\begin{itemize}
    \item \textbf{Network Design:}  
    Min-cut helps in improving the efficiency of communication and optimizing network flow. The algorithm is used in network design to find the minimum capacity cut.

    \item \textbf{Communication Networks:}  
    For understanding the vulnerability of networks to failures, minimum cut can be useful. It helps build robust and fault-tolerant communication networks.

    \item \textbf{VLSI Design:}  
    In Very Large Scale Integration (VLSI) design, the algorithm is useful for partitioning circuits into smaller components, leading to reduced interconnectivity complexity.
\end{itemize}

Reference is made to Section 1.5: \textit{Application: A Randomized Min-Cut Algorithm} in \textit{Probability and Computing (2nd Edition)}.

\hrule
\vspace{0.5cm}

\subsection*{1.2 What is Min-Cut?}

A \textbf{cut-set} in a graph is a set of edges whose removal breaks the graph into two or more connected components.

Given a graph
\[
G = (V, E)
\]
with \( n \) vertices, the \textbf{minimum cut (min-cut) problem} is to find a \textbf{minimum cardinality cut-set} in \( G \).

Min-cut algorithms, such as \textbf{Karger’s algorithm}, are random and can be sensitive to the initial choice of edges. If the algorithm happens to contract critical edges early in the process, it might find a smaller cut.

\subsection*{1.3 Edge Contraction}

The main operation in the algorithm is \textbf{edge contraction}.

Edge contraction is an operation that removes an edge from a graph while simultaneously merging the two vertices.

In contracting an edge \( (u, v) \):

\begin{enumerate}
    \item Merge the two vertices \( u \) and \( v \) into one vertex.
    \item Eliminate all edges connecting \( u \) and \( v \).
    \item Retain all other edges in the graph.
    \item The new graph may have parallel edges but no self-loops.
\end{enumerate}

\hrule
\vspace{0.5cm}

\section*{2. Successful and Unsuccessful Min-Cut Runs}

\subsection*{2.1 Successful Min-Cut Run}

A \textbf{successful min-cut run} refers to the success in the outcome of an algorithm designed to find the minimum cut in a graph.

The figure in the lecture illustrates a sequence of edge contractions that results in the correct minimum cut (Figure: Successful Run).

\subsection*{2.2 Unsuccessful Min-Cut Run}

An \textbf{unsuccessful min-cut run} refers to an iteration of a min-cut algorithm where the algorithm fails to correctly identify the minimum cut of a given graph.

The figure in the lecture illustrates contractions that lead to a non-minimum cut (Figure: Unsuccessful Run).

\hrule
\vspace{0.5cm}

\section*{3. Max-Flow Min-Cut Theorem}

The \textbf{Max-Flow Min-Cut Theorem} states:

\begin{quote}
“In a flow network, the maximum amount of flow passing from the source to the sink is equal to the total weight of the edges in a minimum cut.”
\end{quote}

Definitions:

\begin{itemize}
    \item \textbf{Capacity of a cut:}  
    The sum of the capacity of the edges in the cut that are oriented from a vertex \( \in X \) to a vertex \( \in Y \).

    \item \textbf{Minimum cut:}  
    The cut in the network that has the smallest possible capacity.

    \item \textbf{Minimum cut capacity:}  
    The capacity of the minimum cut.

    \item \textbf{Maximum flow:}  
    The largest possible flow from source \( S \) to sink \( T \).
\end{itemize}

\hrule
\vspace{0.5cm}

\section*{4. Deterministic Min-Cut Algorithm}

\subsection*{4.1 Stoer–Wagner Min-Cut Algorithm}

Let \( s \) and \( t \) be two vertices of a graph \( G \). Let \( G/{s,t} \) be the graph obtained by merging \( s \) and \( t \).

Then a minimum cut of \( G \) can be obtained by taking the smaller of:

\begin{enumerate}
    \item A minimum \( s\text{-}t \)-cut of \( G \), and
    \item A minimum cut of \( G/{s,t} \).
\end{enumerate}

The theorem holds since:

\begin{itemize}
    \item Either there is a minimum cut of \( G \) that separates \( s \) and \( t \), then a minimum \( s\text{-}t \)-cut of \( G \) is a minimum cut of \( G \);
    \item Or there is none, then a minimum cut of \( G/{s,t} \) does the job.
\end{itemize}

\subsection*{4.2 Pseudocode}

\textbf{Algorithm 1: MinimumCutPhase(G, a)}

\begin{enumerate}
    \item \( A \leftarrow \{a\} \)
    \item while \( A \neq V \) do
    \begin{itemize}
        \item add to \( A \) the most tightly connected vertex
    \end{itemize}
    \item return the cut weight as the “cut of the phase”
\end{enumerate}

\textbf{Algorithm 2: MinimumCut(G)}

\begin{enumerate}
    \item while \( |V| \ge 1 \) do
    \begin{itemize}
        \item choose any \( a \) from \( V \);
        \item MinimumCutPhase(G, a);
        \item if the cut-of-the-phase is lighter than the current minimum cut then
        \begin{itemize}
            \item store the cut-of-the-phase as the current minimum cut;
        \end{itemize}
        \item shrink \( G \) by merging the two vertices added last;
    \end{itemize}
    \item return the minimum cut
\end{enumerate}

\hrule
\vspace{0.5cm}

\section*{5. Randomized Min-Cut Algorithm}

\subsection*{5.1 Why Randomized Algorithm?}

Randomized algorithms provide a probabilistic guarantee of success. They provide a more accurate estimate of the minimum cut with fewer iterations.

Characteristics listed in lecture:

\begin{itemize}
    \item Efficiency
    \item Parallelization
    \item Approximation Guarantees
    \item Avoidance of Worst-Case Instances
    \item Heuristic Nature
    \item Robustness
\end{itemize}

\subsection*{5.2 Karger’s Randomized Algorithm}

The lecture provides an example run where random edges are picked in such a way that the output cut is not minimal. The output cut produced is \{a, c, e\}, but the minimal cut is either \{b, e\} or \{a, d\}.

This illustrates the possibility of unsuccessful runs due to random contraction choices.

\subsection*{5.3 Pseudocode: Recursive-Randomized-Min-Cut(G, $\alpha$)}

\textbf{Input:}  
An undirected multigraph \( G \) with \( n \) vertices and an integer constant \( \alpha > 0 \).

\textbf{Output:}  
A cut \( C \) of \( G \).

If \( n \le \alpha^3 \) then:

\begin{itemize}
    \item \( C \leftarrow \) a min-cut of \( G \) found using brute force (exhaustive) search;
\end{itemize}

Else:

\begin{itemize}
    \item for \( i \leftarrow 1 \) to \( \alpha \) do
    \begin{itemize}
        \item \( G' \leftarrow \) multigraph obtained by applying \( n - \frac{n}{\sqrt{\alpha}} \) random contraction steps on \( G \);
        \item \( C' \leftarrow \) Recursive-Randomized-Min-Cut(\( G', \alpha \));
        \item if \( i = 1 \) or \( |C'| < |C| \) then
        \begin{itemize}
            \item \( C \leftarrow C' \);
        \end{itemize}
    \end{itemize}
\end{itemize}

Return \( C \).

\hrule
\vspace{0.5cm}

\section*{6. Comparison: Deterministic vs Randomized Min-Cut}

Any specific problem is responsible for deciding the approach.

\subsection*{Deterministic Min Cut:}

\begin{itemize}
    \item Always guarantees an exact minimum cut.
    \item May have higher time complexity for large graphs.
    \item Stoer–Wagner algorithm has time complexity
    \[
    O(V \cdot E + V^2 \log V)
    \]
    where \( V \) is number of vertices and \( E \) is number of edges.
\end{itemize}

\subsection*{Randomized Min Cut:}

\begin{itemize}
    \item Provides an approximate minimum cut with high probability.
    \item Karger’s algorithm has time complexity
    \[
    O(V^2)
    \]
    where \( V \) is the number of vertices.
\end{itemize}

\hrule
\vspace{0.5cm}

\section*{7. Theorem for Min-Cut Set}

The algorithm outputs a min-cut set with probability at least
\[
\frac{2}{n(n-1)}.
\]

\hrule
\vspace{0.5cm}

\section*{8. Python Simulation}

\begin{itemize}
    \item Open Campuswire Post regarding today’s lecture (L10).
    \item Download the `.ipynb` file pasted in the comments.
\end{itemize}

\bigskip
\centerline{\textbf{End of Lecture 10}}

\end{document}
